% Exact LaTeX recreation of the supplied Phase-1.pdf
% Compile with XeLaTeX.
% Place member pictures in the same folder as this .tex file.

\documentclass[11pt,a4paper]{article}

% -------------------------------------------------
% IMAGE SETTINGS
% -------------------------------------------------
\usepackage[final]{graphicx}
\setkeys{Gin}{draft=false}

% -------------------------------------------------
% PACKAGES
% -------------------------------------------------
\usepackage[a4paper,left=2.0cm,right=2.0cm,top=1.55cm,bottom=1.55cm]{geometry}

% ACTUAL TIMES NEW ROMAN
\usepackage{fontspec}
\setmainfont{Times New Roman}

\usepackage{xcolor}
\usepackage{array}
\usepackage{tabularx}
\usepackage{ragged2e}
\usepackage{enumitem}
\usepackage{fancyhdr}
\usepackage{tikz}
\usepackage{setspace}

% -------------------------------------------------
% COLOURS
% -------------------------------------------------
\definecolor{TitleBlue}{RGB}{61,103,154}

% -------------------------------------------------
% GENERAL SETTINGS
% -------------------------------------------------
\setlength{\parindent}{0pt}
\setlength{\parskip}{0pt}
\renewcommand{\arraystretch}{1.0}

% Make all existing \textit commands normal text
\renewcommand{\textit}[1]{#1}

% -------------------------------------------------
% HEADINGS
% -------------------------------------------------

\newcommand{\blueheading}[1]{%
    {\rmfamily\color{TitleBlue}\fontsize{15}{18}\selectfont #1}%
}

\newcommand{\bluebigheading}[1]{%
    {\rmfamily\bfseries\color{TitleBlue}\fontsize{16}{19}\selectfont #1}%
}

\newcommand{\instruction}[1]{%
    {\rmfamily\fontsize{10.5}{12.5}\selectfont #1}%
}

\newcommand{\answerbox}[2][3.0cm]{%
    \vspace{2pt}
    \noindent
    \fbox{%
        \begin{minipage}[t][#1][t]{\dimexpr\textwidth-2\fboxsep-2\fboxrule\relax}
            \vspace{1mm}
            \rmfamily #2
        \end{minipage}
    }
}

% -------------------------------------------------
% HEADER / FOOTER
% -------------------------------------------------
\pagestyle{fancy}
\fancyhf{}
\renewcommand{\headrulewidth}{0pt}
\renewcommand{\footrulewidth}{0pt}

\fancyhead[L]{%
    \rmfamily\bfseries\fontsize{10.5}{12}\selectfont
     DSA PROJECT BASED LEARNING WITH C++
}

\fancyhead[R]{%
    \rmfamily\bfseries\fontsize{10.5}{12}\selectfont
    PROJECT PROPOSAL
}

\fancyfoot[R]{%
    \rmfamily\fontsize{10}{12}\selectfont
    Page \thepage
}

\setlength{\headheight}{14pt}
\setlength{\headsep}{23pt}
\setlength{\footskip}{24pt}

% -------------------------------------------------
% DOCUMENT
% -------------------------------------------------
\begin{document}

% -------------------------------------------------
% PAGE 1
% -------------------------------------------------

\vspace*{4mm}

\bluebigheading{PROJECT AND TEAM INFORMATION}

\vspace{13mm}

\blueheading{Project Title}

\vspace{1mm}



\answerbox[1.0cm]{VidOptix — VIDEO COMPRESSION SYSTEM}

\vspace{13mm}

\blueheading{Student / Team Information}

\vspace{2mm}

% -------------------------------------------------
% TEAM INFORMATION TABLE
% -------------------------------------------------

\noindent
\begin{tabularx}{\textwidth}
{|>{\RaggedRight\arraybackslash}p{0.69\textwidth}|
 >{\RaggedRight\arraybackslash}X|}
\hline

\begin{minipage}[t]{\linewidth}
\vspace{1.5mm}
\textit{Team Name:}\\[-1mm]
\textit{Team \#}
\vspace{1.5mm}
\end{minipage}
&
\begin{minipage}[t]{\linewidth}
\vspace{1.5mm}
\textit{BitSquad}\\[-1mm]
\textit{DSCPP-III-2026-T198}
\vspace{1.5mm}
\end{minipage}
\\
\hline

% -------------------------------------------------
% TEAM MEMBER 1
% -------------------------------------------------

\begin{minipage}[t][3.25cm][t]{\linewidth}
\vspace{1.5mm}
\textbf{Team member 1 (Team Lead)}\\[-1mm]
\vspace{2.0cm}
\end{minipage}
&
\begin{minipage}[t][3.25cm][t]{\linewidth}
\vspace{1.5mm}

\textit{Goswami, Anshika -- 2510380062}\\
\textit{2510380062@geu.ac.in}

\vspace{2mm}

\includegraphics[
    width=3.0cm,
    height=1.5cm,
    keepaspectratio
]{WhatsApp Image 2026-08-28 at 8.21.27 PM.jpeg}

\end{minipage}
\\
\hline

% -------------------------------------------------
% TEAM MEMBER 2
% -------------------------------------------------

\begin{minipage}[t][3.25cm][t]{\linewidth}
\vspace{1.5mm}
\textbf{Team member 2}\\[-1mm]
\vspace{2.0cm}
\end{minipage}
&
\begin{minipage}[t][3.25cm][t]{\linewidth}
\vspace{1.5mm}

\textit{ Yadav ,Yaduveer-- 2510370438}\\
\textit{2510370438@geu.ac.in}

\vspace{2mm}

\includegraphics[
    width=3.0cm,
    height=1.5cm,
    keepaspectratio
]{WhatsApp Image 2026-08-28 at 7.54.24 PM.jpeg}

\end{minipage}
\\
\hline

% -------------------------------------------------
% TEAM MEMBER 3
% -------------------------------------------------

\begin{minipage}[t][3.25cm][t]{\linewidth}
\vspace{1.5mm}
\textbf{Team member 3}\\[-1mm]

\vspace{2.0cm}
\end{minipage}
&
\begin{minipage}[t][3.25cm][t]{\linewidth}
\vspace{1.5mm}

\textit{Chaudhary, Deepanshi -- 2510370456}\\
\textit{2510370456@geu.ac.in}

\vspace{2mm}

\includegraphics[
    width=3.0cm,
    height=1.5cm,
    keepaspectratio
]{WhatsApp Image 2026-08-28 at 8.36.05 PM.jpeg}

\end{minipage}
\\
\hline

% -------------------------------------------------
% TEAM MEMBER 4
% -------------------------------------------------

\begin{minipage}[t][3.25cm][t]{\linewidth}
\vspace{1.5mm}
\textbf{Team member 4}\\[-1mm]

\vspace{2.0cm}
\end{minipage}
&
\begin{minipage}[t][3.25cm][t]{\linewidth}
\vspace{1.5mm}

\textit{Ubaid,Mohd -- 2510370084}\\
\textit{2510370084@geu.ac.in}

\vspace{2mm}

\includegraphics[
    width=3.0cm,
    height=1.5cm,
    keepaspectratio
]{WhatsApp Image 2026-08-28 at 8.50.58 PM.jpeg}

\end{minipage}
\\
\hline

\end{tabularx}

\newpage

% -------------------------------------------------
% PAGE 2
% -------------------------------------------------
\vspace{8mm}
\bluebigheading{PROPOSAL DESCRIPTION (10 pts)}

\vspace{12mm}

\blueheading{Motivation (1 pt)}

\vspace{3mm}


\answerbox[7.5cm]{
\textbf{Problem Statement}

\vspace{2mm}

Videos contain a very large amount of data, requiring significant \textbf{storage space and bandwidth} for saving, sharing, and transferring. As video resolution and duration increase, their file sizes also grow, making efficient compression important.

Our project aims to develop a \textbf{simple and efficient video compression system} that reduces video file size while maintaining approximately \textbf{75--80\% of the original visual quality}.

\vspace{1mm}

The main problem we aim to solve is the \textbf{trade-off between video size and visual quality}. Instead of preserving every detail and producing a large file, the system will allow controlled quality loss to achieve meaningful size reduction.

The project will also demonstrate the practical application of data structures and algorithms, particularly Huffman coding, priority queues/min-heaps, and tree structures, in a real-world compression problem.
}

\vspace{12mm}

\blueheading{State of the Art / Current solution (1 pt)}

\vspace{1mm}

\answerbox[7.95cm]{Today, video compression is mainly performed using standardized codecs such as H.264/AVC, H.265/HEVC, and AV1. These codecs combine several techniques to reduce video size while maintaining visual quality. Modern video codecs do not only compress individual frames; they also exploit similarities between consecutive frames using motion estimation and inter-frame prediction.

Common techniques used in current video compression include 

\vspace{1mm}

\textbf{Transform coding:} Converts image data into frequency components to identify less-important information.

\textbf{Quantization:} Reduces less-important information to achieve controlled quality loss.

\textbf{Zigzag Scanning:} Arranges coefficients to improve subsequent compression.

\textbf{RLE:} Compresses repeated values, particularly zeros.

\textbf{Huffman Coding:} Further reduces the data size using variable-length encoding.

\vspace{2mm}
Modern codecs provide excellent compression performance, but their implementations are complex and involve many interacting stages.
\vspace{1mm}

Our project uses a simplified frame-based approach consisting of DCT, quantization, zigzag scanning, RLE, and Huffman coding. It focuses on understanding and implementing the fundamental compression techniques rather than reproducing the complexity of modern codecs. The goal is to achieve meaningful size reduction while maintaining approximately 75–80 percent visual quality.
}

\vspace{25mm}
\newpage

\blueheading{Project Goals and Milestones (2 pts)}

\vspace{3mm}

\answerbox[8.95cm]{\textbf{Project Goals and Milestones}

\begin{tabular}{|p{0.38\linewidth}|p{0.52\linewidth}|}
\hline
\textbf{Goals} & \textbf{Milestones} \\
\hline

Develop a simple video compression system that reduces video size while maintaining approximately 75--80\% visual quality.
& Problem identification, topic selection, and initial research on video compression -- \textbf{Completed} \\
\hline

Design and implement a manageable compression pipeline using DCT, quantization, zigzag scanning, RLE, and Huffman coding.
& Finalization of the compression pipeline and selection of OpenCV/FFmpeg for video handling -- \textbf{Completed} \\
\hline

Implement and integrate the proposed compression techniques into a functional system.
& Implementation and integration of DCT, quantization, zigzag, RLE, and Huffman coding -- \textbf{Week 1--5} \\
\hline

Evaluate and optimize the system for a suitable balance between compression and visual quality.
& Testing and evaluation using compression ratio, file size, processing time, and PSNR/SSIM -- \textbf{Week 6--8} \\
\hline

& Final optimization, documentation, and project demonstration -- \textbf{Week 9--10} \\
\hline

\end{tabular}}

\vspace{12mm}

\blueheading{Project Approach (3 pts)}

\vspace{3mm}


\answerbox[10.05cm]{The project will follow a modular, frame-based video compression approach. The solution will be designed as a sequence of compression stages, where each stage processes the output of the previous stage.


\vspace{3mm}

\textbf{1.Video Processing and Frame Extraction:}
The input video will be read using OpenCV/FFmpeg, which will handle video decoding, frame extraction, and basic properties such as resolution and frame rate. Processing will be done frame-by-frame to keep the system simple and memory-efficient.

\textbf{2.Transform and Lossy Compression:}
Each frame will be divided into 8×8 pixel blocks. DCT will convert the pixel values into frequency coefficients, followed by quantization to reduce the precision of less-important coefficients. The quantization level will be adjusted to control the balance between compression and visual quality.

\textbf{3.Coefficient Processing and RLE:}
The quantized coefficients will be arranged using zigzag scanning, which places similar and zero-valued coefficients closer together. Run-Length Encoding (RLE) will then represent consecutive repeated values more efficiently.

\textbf{4.Huffman Coding and Compressed Data:}
The RLE output will be further compressed using Huffman coding. A priority queue/min-heap and Huffman tree will be implemented in C/C++ to generate variable-length codes. The resulting compressed binary data will be stored along with the required metadata.

\textbf{5.Testing and Optimization:}
Different quantization levels and sample videos will be tested to achieve a suitable balance between file-size reduction and approximately 75–80 percent visual quality. The system will be evaluated using compression ratio, file-size reduction, processing time, and PSNR/SSIM. 
}
\newpage

% -------------------------------------------------
% PAGE 3
% -------------------------------------------------
\blueheading{System Architecture (High Level Diagram)(2 pts)}

\vspace{1mm}

\answerbox[10.8cm]{
    \centering
    \includegraphics[
        width=\linewidth,
        keepaspectratio
    ]{architecture.png}
}

\vspace{12mm}


\blueheading{Project Outcome / Deliverables (1 pts)}

\vspace{1mm}


\answerbox[8.5cm]{The project will result in a functional frame-based video compression system that reduces video file size while maintaining approximately 75–80 percent visual quality.
\vspace{2mm}

Expected outcomes and deliverables:

\textbf{1.}A working video compression pipeline using DCT, quantization, zigzag scanning, RLE, and Huffman coding.

\vspace{1mm}

\textbf{2.}A program capable of taking a video as input, processing its frames, and generating a compressed data output.

\vspace{1mm}

\textbf{3.}A structured implementation of Huffman coding using a priority queue/min-heap as the primary DSA component.

\vspace{1mm}

\textbf{4.}A quality and compression evaluation based on compression ratio, file-size reduction, processing time, and PSNR/SSIM.

\vspace{1mm}

\textbf{5.}A comparison of different quantization levels to identify a suitable compression–quality trade-off.

\vspace{1mm}

\textbf{6.}Complete source code, system architecture/design, documentation, test results, and final project demonstration.

\vspace{3mm}
The final deliverable will demonstrate that the proposed combination of transform, statistical, and run-length compression techniques can provide meaningful video-size reduction while retaining acceptable visual quality.}

\vspace{12mm}
\newpage

\blueheading{Assumptions}

\vspace{6mm}

\answerbox[2.65cm]{Key assumptions include:

Approximately 75–80percent visual quality is considered acceptable.
Videos used for testing will have moderate resolution and duration to keep processing time manageable.
OpenCV/FFmpeg will handle video reading and frame extraction.
The proposed DCT, quantization, zigzag, RLE, and Huffman stages will handle the core compression process.}

\vspace{12mm}

\blueheading{References}

\vspace{1mm}


\answerbox[6.8cm]{
\begin{enumerate}
\setlength{\itemsep}{0pt}
\setlength{\parskip}{0pt}

\item OpenCV Documentation -- Video I/O and Frame Processing: \\
\href{https://docs.opencv.org/4.13.0/dd/de7/group}
\vspace{2mm}
\item OpenCV -- Reading and Writing Videos using OpenCV: \\
\href{https://opencv.org/reading-and-writing-videos-using-opencv/}{https://opencv.org/reading-and-writing-videos-using-opencv/}
\vspace{2mm}
\item NIST Dictionary of Algorithms and Data Structures -- Huffman Coding: \\
\href{https://xlinux.nist.gov/dads/HTML/huffmanCoding.html}
\vspace{2mm}
\item Wolfram MathWorld -- Huffman Coding: \\
\href{https://mathworld.wolfram.com/HuffmanCoding.html}
\vspace{2mm}
\item RLE Discussion and Implementation: \\
\href{https://michaeldipperstein.github.io/rle.html}
\vspace{2mm}

\end{enumerate}
}

\end{document}